% ============================================================
%  style/layout.tex — Pakete, Seitenränder, Farben, Hintergrund, Kopf/Fuß, Inhaltsverzeichnis
%  (wird von style.tex eingebunden)
% ============================================================
% ---------- Geometrie für Zweiseitigen Druck ----------
\usepackage[
  a5paper,
  twoside,
  inner=14mm,
  outer=10mm,
  top=12mm,
  bottom=10mm,
  footskip= 0mm,
  bindingoffset=4mm
]{geometry}

\usepackage[utf8]{inputenc}
\usepackage[T1]{fontenc}
\usepackage{xcolor}
\usepackage{alltt}          % Erlaubt LaTeX-Befehle in vorformatiertem Text
\usepackage{fancyhdr}
\usepackage{tikz}
\usetikzlibrary{calc}
\usepackage{eso-pic}
\usepackage[colorlinks=false,pdfborder={0 0 0}]{hyperref}
\usepackage{titlesec}
\usepackage{tocloft}

% ---------- Farbpalette ----------------------------------------
\definecolor{paperbg}{HTML}{fbfbf9}   % Papierfarbe (Vintage-Creme)
\definecolor{inkmain}{HTML}{000000}   % Haupttext
\definecolor{inkchords}{HTML}{232323} % Akkorde
\definecolor{inkheading}{HTML}{282828}% Überschriften
\definecolor{inkmeta}{HTML}{8e8e8e}   % Interpret / Capo-Zeile
\definecolor{inksketch}{HTML}{b8b6b6} % Skizzen / Griffbilder
\definecolor{rulecolor}{HTML}{999999} % Trennlinien / Punktlinien

% ---------- Seitenhintergrund ----------
\AddToShipoutPictureBG{%
  \AtPageLowerLeft{%
    \begin{tikzpicture}[x=1pt,y=1pt]
      \fill[paperbg] (0,0) rectangle (\paperwidth,\paperheight);
    \end{tikzpicture}%
  }%
}

\color{inkmain}

% ---------- Seitenstil & Seitenzahlen (Zweiseitig) ----------
\pagestyle{fancy}
\fancyhf{}
\renewcommand{\headrulewidth}{0pt}
\renewcommand{\footrulewidth}{0pt}

\fancyfoot[RO,LE]{\serif\fontsize{\SzPageNum}{\LdPageNum}\selectfont\color{inkmeta}\thepage}

\fancypagestyle{plain}{%
  \fancyhf{}%
  \renewcommand{\headrulewidth}{0pt}%
  \renewcommand{\footrulewidth}{0pt}%
  \fancyfoot[RO,LE]{\serif\fontsize{\SzPageNum}{\LdPageNum}\selectfont\color{inkmeta}\thepage}%
}

% ---------- Schriften ----------
\newcommand{\songnumfont}{\serif\bfseries\fontsize{\SzSongNum}{\LdSongNum}\selectfont\color{inkheading}}
\newcommand{\songtitlefont}{\serif\bfseries\fontsize{\SzSongTitle}{\LdSongTitle}\selectfont\color{inkheading}}
\newcommand{\songmetafont}{\serif\itshape\fontsize{\SzMeta}{\LdMeta}\selectfont\color{inkmeta}}
\newcommand{\songbodysize}{\fontsize{\SzBody}{\LdBody}\selectfont}
\newcommand{\covertitlefont}{\serif\bfseries\fontsize{\SzCoverTitle}{\LdCoverTitle}\selectfont\color{inkheading}}
\newcommand{\coversubfont}{\serif\itshape\fontsize{\SzCoverSub}{\LdCoverSub}\selectfont\color{inkmeta}}
\newcommand{\cattitlefont}{\serif\bfseries\fontsize{\SzCat}{\LdCat}\selectfont\color{inkheading}}

% ---------- Inhaltsverzeichnis-Optik ----------
\renewcommand{\contentsname}{Inhaltsverzeichnis}
\renewcommand{\cftsecfont}{\serif\color{inkmain}}
\renewcommand{\cftsecpagefont}{\serif\color{inkmeta}}
\renewcommand{\cftsecleader}{\cftdotfill{\cftdotsep}}
\renewcommand{\cftdot}{\textcolor{rulecolor}{.}}
\setlength{\cftbeforesecskip}{1.1mm}
\renewcommand{\cftsecafterpnum}{\vskip 0pt}

\renewcommand{\cftpartfont}{\serif\bfseries\color{inkheading}}
\renewcommand{\cftpartpagefont}{\serif\color{inkheading}}
\setlength{\cftbeforepartskip}{4mm}
\renewcommand{\cftpartleader}{\hfill}

% Feste Breite für Nummer/Buchstabe im Inhaltsverzeichnis,
% damit alle Titel gleich eingerückt beginnen (auch ab Nr. 10)
\newlength{\tocnumw}\setlength{\tocnumw}{6mm}

